\section{The Realizability Barrier}

A natural hope for a finite tractability taxonomy would be this: perhaps optimizer-induced quotients realize only a narrow subclass of equivalence relations, so the finite frontier would emerge from realizability alone. In the unconstrained setting, that hope is false.

\leanmetapending{\LHrng{FR}{7}{8}}
\begin{theorem}[Every Labeling Kernel Is Optimizer-Realizable]\label{thm:labeling-realizable}
Let $\phi : S \to T$ be any function. There exists a decision problem with action space $T$ and state space $S$ such that for every state $s \in S$,
\[
\Opt(s) = \{\phi(s)\}.
\]
Consequently, for all states $s,s' \in S$,
\[
\Opt(s)=\Opt(s') \quad\Longleftrightarrow\quad \phi(s)=\phi(s').
\]
Equivalently, the decision quotient of the constructed problem is exactly the kernel partition of $\phi$.
\end{theorem}

\begin{proof}
Take the action set to be $T$ itself and define
\[
U(a,s)=
\begin{cases}
1 & \text{if } a=\phi(s),\\
0 & \text{otherwise.}
\end{cases}
\]
Then the unique maximizer at state $s$ is the designated action $\phi(s)$, so the optimizer set is exactly $\{\phi(s)\}$. Two states therefore have the same optimizer set if and only if they receive the same label under $\phi$.
\end{proof}

Theorem~\ref{thm:labeling-realizable} is the key negative fact for the frontier program. It shows that optimizer realizability by itself is extremely permissive. If arbitrary label kernels are already optimizer-induced, then no finite taxonomy can come merely from asking which abstract quotients are realizable.

\leanmetapending{\LHrng{FR}{45}{47}}
\begin{proposition}[Every Equivalence Relation Is Optimizer-Realizable]\label{prop:setoid-realizable}
For every equivalence relation on the state space, there exists a decision problem whose decision quotient relation is exactly that equivalence relation. Moreover, the quotient object of the realizing problem is canonically equivalent to the corresponding setoid quotient.
\end{proposition}

\begin{proof}
Let $\approx$ be an equivalence relation on $S$, and let $\pi : S \to S/{\approx}$ be the quotient map. Apply Theorem~\ref{thm:labeling-realizable} to the labeling map $\pi$. The resulting decision problem satisfies
\[
\Opt(s)=\Opt(s') \iff \pi(s)=\pi(s'),
\]
and the right-hand side is exactly the original equivalence relation. Since the realizing quotient identifies states precisely by equality of quotient labels, its quotient object is canonically equivalent to the setoid quotient itself.
\end{proof}

Proposition~\ref{prop:setoid-realizable} is the maximal realizability statement available in the present framework. The labeling-kernel theorem is a special case. Thus the realizability barrier is not merely broad; at the level of abstract quotient relations, it is maximal.

\leanmetapending{\LHrng{FR}{39}{40}}
\begin{proposition}[The Realized Quotient Is Exactly the Label Range]\label{prop:realized-quotient-range}
For every labeling map $\phi : S \to T$, the decision quotient of the realizing problem is canonically equivalent to the range of $\phi$. In finite settings, the quotient cardinality is therefore exactly $|\operatorname{range}(\phi)|$.
\end{proposition}

\begin{proof}
By Theorem~\ref{thm:labeling-realizable}, two states lie in the same decision class if and only if they receive the same label under $\phi$. Therefore the quotient class of a state depends only on the value $\phi(s)$, and sending the class of $s$ to $\phi(s)$ defines a well-defined map from the decision quotient to $\operatorname{range}(\phi)$. This map is surjective by definition of the range and injective because equality in the range means equality of labels, hence equality of quotient classes.
\end{proof}

Proposition~\ref{prop:realized-quotient-range} strengthens Theorem~\ref{thm:labeling-realizable}. The realizing construction does not merely realize the induced equivalence relation abstractly; it realizes the quotient object itself with exactly the expected image size.

\leanmetapending{\LHrng{FR}{7}{8}}
\begin{corollary}[Realizability Alone Cannot Be the Frontier]\label{cor:realizability-alone-not-enough}
Any future optimizer-compatible tractability metatheorem must use more than quotient realizability alone. Additional restrictions have to enter through representation, coordinate structure, utility structure, closure properties of the family under study, or comparable algebraic invariants.
\end{corollary}

This shifts the burden of the metatheorem. The hard question is not which quotients can arise in principle; almost all kernels can. The hard question is which \emph{natural structural families of decision problems} force tractability or preserve hardness once coordinate structure and utility representation are taken seriously.
